\documentclass[11pt, a4paper]{scrartcl}
% Packages
\usepackage[margin=1.25in]{geometry}
\usepackage{index}
\usepackage{amsbsy} % Bold math symbols
\makeindex
\usepackage{tcolorbox}
\usepackage{varwidth}
\usepackage{amsmath, amssymb}
\usepackage{esint}
\usepackage{titlesec}
\usepackage{xcolor}
\usepackage{titling}
\usepackage[linktocpage]{hyperref}
\usepackage{pgfplots}
\usepackage{caption}
\usepackage{amsthm}
\usepackage{import}
\usepackage{cancel}
\usepackage{caption}
\usepackage{nicematrix}
\usepackage{mathtools}
\usepackage{enumerate}
\usepackage{graphicx}
\usepackage{makecell}
\usepackage{tikz}
\usepackage[version=4]{mhchem}
\usetikzlibrary{arrows.meta, calc ,positioning, angles,quotes,patterns.meta,intersections}
\usepackage[italian]{babel}
% Font
%\usepackage{fontspec}
%\setmainfont{Libertinus Serif}

%\usepackage{unicode-math}
%\setmathfont{Libertinus Math}


% To reset footnote numbering each page
\usepackage[perpage]{footmisc}
\usepackage{setspace}
\setstretch{1.2}

% Titles 
\title{Note di\\ \vspace{.15cm} Fisica nucleare}
\author{}
\date{}


% svolgimento
\newenvironment{svolgimento}{\renewcommand\qedsymbol{$\blacksquare$}\begin{proof}[Svolgimento]}{\end{proof}}

\newtheoremstyle{style}% name of the style to be used
{5pt}% measure of space to leave above the theorem. E.g.: 3pt
{5pt}% measure of space to leave below the theorem. E.g.: 3pt
{\normalfont}% name of font to use in the body of the theorem
%{15pt}% measure of space to indent
{0pt}% measure of space to indent
{\noindent\bfseries}% name of head font
{}% punctuation between head and body
{ }% space after theorem head; " " = normal interword space
{\thmname{#1}\thmnumber{ #2}{\thmnote{ (#3)}.\ }}


\theoremstyle{style}
\newtheorem{esempio}{Esempio}[section]
\newtheorem{definizione}{Definizione}[section]
\newtheorem{prop}{Proposizione}[section]
\newtheorem{teorema}{Teorema}[section]
\newtheorem{lemma}{Lemma}[teorema]
\newtheorem{corollario}{Corollario}[teorema]
\newtheorem{osservazione}{Osservazione}[section]
\newtheorem{assunzione}{Assunzione}[section]
\newtheorem{notazione}{Notazione}[section]

\newtheoremstyle{style1}% name of the style to be used
{5pt}% measure of space to leave above the theorem. E.g.: 3pt
{5pt}% measure of space to leave below the theorem. E.g.: 3pt
{\normalfont}% name of font to use in the body of the theorem
%{15pt}% measure of space to indent
{0pt}% measure of space to indent
{\noindent\bfseries\color{red}}% name of head font
{}% punctuation between head and body
{ }% space after theorem head; " " = normal interword space
{\thmname{#1}\thmnumber{ #2}{\thmnote{ (#3)}.\ }}

\theoremstyle{style1}
\newtheorem{esercizio}{Esercizio}[section]


%% Generic box
\newtcolorbox{eqbox}[1][]
{
colback=gray!10,
arc=0pt,
boxrule=0pt,
title=#1
}

 \newenvironment{boxenv}[1][]{
    \begin{eqbox}[#1]
    }{
   \end{eqbox}
}

%Captions
\captionsetup[figure]{font=footnotesize,labelfont=footnotesize}
\captionsetup[table]{font=footnotesize,labelfont=footnotesize}
%Titlesec
\titleformat{\section}
{\fontsize{15}{20}\sffamily\scshape}
{\normalfont\color{gray}{\fontsize{20}{20}\selectfont\thesection}}
{0.7em}
{}
\hypersetup{colorlinks,breaklinks, linkcolor=[RGB]{74, 122, 164}}
\definecolor{asdf}{HTML}{4a7aa4}
% Personalizza la formattazione della subsection
\titleformat{\subsection}[block]{\fontsize{13}{20}\bfseries}{\normalfont\thesubsection}{.5em}{}


% Personalizza la formattazione della subsubsection
\titleformat{\subsubsection}[block]{\fontsize{12}{20}\bfseries}{\normalfont\thesubsubsection}{.5em}{}

% Maketitle customization
\renewcommand{\maketitle}{
\begin{center}
{\sffamily
{\fontsize{20}{20}\selectfont\MakeUppercase\thetitle}}

\vspace{0.2in}

{\large\scshape\sffamily\theauthor}
\end{center}
}


% Differenziale 
\newcommand{\dd}{\mathop{}\!\mathrm{d}}




%Evaluate symbol
\DeclareMathOperator{\di}{d\!}
\newcommand*\Eval[3]{\left.#1\right\rvert_{#2}^{#3}}

%%%%%%% Numero delle equazioni in formato a.b
\numberwithin{equation}{subsection}
%%%%%

%%%%%%%%%% Personalizzazione numeri lista
\renewcommand{\theenumi}{(\arabic{enumi})}

%%%% Table of contents

\usepackage[titles]{tocloft}

\renewcommand{\cftdot}{}
\usepackage{titletoc}
%\setcounter{tocdepth}{2}

%%%%%%%%%%%%%%%% Toc style

% Personalizzazione scritta indice


%%%% Lambda bar
\newcommand{\lbar}{{\mkern0.75mu\mathchar '26\mkern -9.75mu\lambda}}


\begin{document}
\maketitle
\newpage
\tableofcontents 
\newpage
\section{Propriet\`a generali dei nuclei atomici}
Le scale rilevanti per lo studio dei nuclei atomici sono:
\begin{itemize}
	\item \textit{lunghezze} di $\sim 10^{-15}$ m;
	\item \textit{densit\`a} di $\sim 10^{14}\text{ g}/\text{cm}^3$;
	\item \textit{energie} di $\sim 1$ MeV.
\end{itemize}
Le forze rilevanti, invece, sono dovute all'interazione nucleare forte e debole.
\subsection{Modello nucleare p-e}
Si assume che il nucleo atomico sia composto da elettroni e protoni distribuiti omogeneamente a formare una sfera di raggio $R$.
In questo modello, si indica con $A$ il numero dei protoni e con $A-Z$ il numero degli elettroni nel nucleo.
Allora, la massa del nucleo \`e data da 
\[
M = A\cdot m_p + (A-Z) \cdot m_e \approx A\cdot m_p
\] 
essendo $m_p \gg m_e$.
Questo \`e in accordo con quanto teorizzato da Aston (1922), per cui $M = A \cdot m_{\ce{^1H}}$.
La carica del nucleo, allora, si trova essere 
\[
Q = A\cdot e - (A-Z) \cdot e = Z\cdot e
\] 
Cos\`i facendo, l'atomo \`e globalmente neutro se gli orbitano attorno $Z$ elettroni.
Ciononostante, il modello presenta una serie di problemi.
\begin{itemize}
	\item \textbf{Energia di legame degli elettroni.} 

		Si fissa $R = 5$ Fm e, usando la relazione di indeterminazione di Heisenberg, si trova che 
		\[
		\begin{cases}
			\Delta p_x \Delta x \sim \frac{\hbar }{2}\\
			\Delta x \sim R\\
			\Delta p_x \sim p_x
		\end{cases} \implies p_x \sim \frac{\hbar }{2R} = \frac{1}{2} \frac{\hbar c }{R} \frac{1}{c}
		\] 
		con $\hbar c \approx 200 $ MeV Fm. 
		Allora
		\[
		p c = c \sqrt{p_x^2 + p_y^2 + p_z^2} \sim \frac{\sqrt{3} }{2} \frac{\hbar c}{R} \approx 34.6 \text{ MeV}
		\] 
		Trascurando il termine $m_e c^2$, si conclude che
		\[
		E = \sqrt{(pc)^2 + (m_e c^2)^2} \sim T_e = pc \approx 34.6 \text{ MeV}
		\] 
		con $T_e$ energia cinetica dell'elettrone.
		L'energia potenziale di un elettrone sulla superficie di un nucleo interamente composto da protoni \`e ottenibile tramite il teorema di Gauss, secondo cui questo nucleo si pu\`o assimilare ad una carica puntiforme di carica pari alla carica totale dell'insieme di protoni:
		\[
			U_\text{coul} = - \frac{e^2 A}{ 4\pi \varepsilon _0 R} = - \underbracket{\frac{e^2}{4 \pi \varepsilon _0 R \hbar c}}_{= \alpha } \hbar c \frac{A}{R}
		\] 
	con $\alpha \approx 1 / 137$ costante di struttura fine.	
	Per un atomo di media grandezza, quindi con $A = 40$ per esempio, usando $R = 5$ Fm, si trova che $U_\text{coul} \approx -12 $ MeV.
	L'energia di un elettrone si stima essere, quindi, una quantit\`a positiva, per cui il nucleo non rimarrebbe legato.
\item \textbf{Momento angolare totale e spin.} 
	Il momento angolare totale \`e dato da $\mathbf{J} = \mathbf{L}  + \mathbf{S} $, con $\mathbf{L} $ che ha associato un numero quantico intero.
	Quanto allo spin, invece, si ha 
	\[
	\mathbf{S}  = \sum_{i=1}^{A} \mathbf{s} _{p_i} + \sum_{j=1}^{A-Z} \mathbf{s} _{e_j}
	\] 
	cio\`e \`e dato dallo spin totale dei protoni, pi\`u quello totale degli elettroni.
	Prendendo per esempio l'azoto-14, che secondo questo modello \`e composto da 14 protoni e 7 elettroni, risulta $\mathbf{S} $ con numero quantico semi-intero (visto che $14 + 7 = 21$ \`e dispari), quindi anche $\mathbf{J} $ ha numero quantico semi-intero; tuttavia, si \`e misurato sperimentalmente che tale momento angolare totale sarebbe dovuto essere pari a $1$.
\end{itemize}





\end{document}
